@ID: OW19D1P1I
@BIDANG: Analisis Real

Diberikan barisan bilangan real $(x_n)$. Jika

$$
\sum_{n=1}^{\infty}x_n<\infty,
$$

maka

$$
\lim_{n\to\infty}\frac{\sum_{k=1}^{\infty}\sqrt{x_k}}{\sqrt n}=\ldots
$$

@ID: OW19D1P2I
@BIDANG: Analisis Real

Diberikan fungsi kontinu $f$ dan $g$, dengan $f:[0,\infty)\to\mathbb R$ dan $g:[0,1]\to\mathbb R$. Jika

$$
\lim_{x\to\infty}f(x)=A
$$

dan

$$
\int_0^1g(x)\,dx=B,
$$

maka

$$
\lim_{n\to\infty}\frac1n\int_0^n f(x)g\!\left(\frac xn\right)\,dx=\ldots
$$

@ID: OW19D1P3I
@BIDANG: Struktur Aljabar

Misalkan

$$
R=\mathbb Z_5[x]/I,
$$

dengan

$$
I=\langle x^2-1\rangle
$$

adalah ideal yang dibangun oleh polinom $x^2-1$ di $\mathbb Z_5[x]$. Banyaknya homomorfisma gelanggang yang bijektif dari $R$ ke $R$ adalah $\ldots$

@ID: OW19D1P4I
@BIDANG: Struktur Aljabar

Misalkan $(G,\circ)$ adalah grup dengan operasi komposisi fungsi dan $W=\mathbb Z_2\times\mathbb Z_2$ merupakan ruang vektor atas $\mathbb Z_2$. Jika $G$ merupakan himpunan yang unsur-unsurnya berupa pemetaan linear $g:W\to W$, dengan sifat untuk setiap basis $\beta$, himpunan

$$
g(\beta)=\{g(b):b\in\beta\}
$$

membentuk basis $W$, maka banyaknya subgrup berorde $2$ di $G$ adalah $\ldots$

@ID: OW19D1P5I
@BIDANG: Struktur Aljabar

Jika $H$ adalah grup yang dibangun oleh $\alpha$ dan $\beta$, dengan

$$
\alpha^2=\beta^{2019}=(\alpha\beta)^2=1,
$$

maka orde dari $H$ adalah $\ldots$

@ID: OW19D1P6I
@BIDANG: Kombinatorika

Sebuah toko menjual empat jenis kembang gula rasa mangga, jeruk, durian, dan kopi. Untuk keperluan sampel akan dipilih paling banyak $3$ rasa mangga, paling banyak $3$ rasa jeruk, paling banyak $2$ rasa durian, dan paling banyak $2$ rasa kopi. Banyaknya cara untuk memilih sampel berukuran $5$ adalah $\ldots$

@ID: OW19D2P1I
@BIDANG: Aljabar Linear

Matriks $A$ adalah matriks berukuran $n\times n$ dengan entri-entri bilangan asli genap dan ganjil yang berbeda satu dengan yang lain. Agar $A$ menjadi matriks nonsingular, maka banyaknya entri bilangan ganjil paling sedikit adalah $\ldots$

@ID: OW19D2P2I
@BIDANG: Aljabar Linear

Misalkan matriks $A\in\mathbb R^{n\times n}$ dan $I$ adalah matriks identitas berukuran $n\times n$. Misalkan pula matriks

$$
B=
\begin{pmatrix}
A+I & A+2I\\
0 & A+3I
\end{pmatrix}.
$$

Jika $2$ adalah salah satu nilai eigen dari $A$, maka nilai-nilai eigen dari $B$ yang dapat diketahui adalah $\ldots$

@ID: OW19D2P3I
@BIDANG: Analisis Kompleks

Jika

$$
A=\{z\in\mathbb C: z^6=1\}
$$

dan

$$
B=\{z\in\mathbb C: z^9=1\},
$$

maka banyaknya anggota $A\cup B$ adalah $\ldots$

@ID: OW19D2P4I
@BIDANG: Analisis Kompleks

Jika diketahui fungsi

$$
f(z)=f(x+iy)=x^2+axy+by^2+i(cx^2+dxy+y^2)
$$

analitik di seluruh bidang kompleks, maka nilai $f(i)$ adalah $\ldots$

@ID: OW19D2P5I
@BIDANG: Analisis Kompleks

Koefisien suku yang memuat $z^{4039}$ pada ekspansi deret Taylor fungsi

$$
f(z)=z\sinh(z^2)
$$

di $z=0$ adalah $\ldots$

@ID: OW19D2P6I
@BIDANG: Kombinatorika

Sebuah tes terdiri atas $10$ soal. Setiap soal diberi nilai bulat dan paling sedikit diberi nilai $5$. Bila soal pertama hanya boleh diberi nilai $10$ atau $15$ dan total nilai tes adalah $100$, banyaknya cara memberi nilai pada tes tersebut adalah $\ldots$

@ID: OW19D1P1U
@BIDANG: Analisis Real

Diberikan fungsi terdiferensial $g:\mathbb R\to\mathbb R$, terdapat $x_0\in\mathbb R$ dengan $g(x_0)=0$ dan untuk setiap $x\in\mathbb R$ berlaku

$$
g(x)+g'(x)>0.
$$

Buktikan bahwa

$$
[g(x)]^{2019}\ge0
$$

untuk setiap $x\ge x_0$.

@ID: OW19D1P2U
@BIDANG: Analisis Real

Diketahui $\Omega$ adalah koleksi semua fungsi $f:\mathbb R\to\mathbb R$ yang memenuhi

$$
|f(x)-f(t)|\le |x-t|^4
$$

untuk setiap $x,t\in\mathbb R$.

Buktikan bahwa:

(a) untuk setiap $f\in\Omega$, $f$ merupakan fungsi terbatas pada $\mathbb R$.

(b) terdapat barisan fungsi $(f_n)$ di dalam $\Omega$ dengan sifat untuk setiap $x\in\mathbb R$ berlaku

$$
\lim_{n\to\infty}f_n(x)=\infty.
$$

@ID: OW19D1P3U
@BIDANG: Struktur Aljabar

Definisi: Dua bilangan bulat $a$ dan $b$ disebut membangun gelanggang $\mathbb Z$ jika untuk setiap $c\in\mathbb Z$ berlaku

$$
c=va+wb
$$

untuk suatu $v,w\in\mathbb Z$.

Misalkan $P$ merupakan peluang dua bilangan bulat yang dipilih secara acak dapat membangun $\mathbb Z$. Buktikan bahwa

$$
P=\frac6{\pi^2}.
$$

(Petunjuk: gunakan deret $\sum_{n=1}^{\infty}\frac1{n^2}$.)

@ID: OW19D1P4U
@BIDANG: Kombinatorika

Misalkan $n$ merupakan bilangan bulat positif yang tidak habis dibagi oleh $2$ maupun $5$. Perlihatkan bahwa terdapat bilangan bulat $q$ yang merupakan kelipatan dari $n$ sedemikian sehingga semua digit dari $q$ adalah $1$.

(Contoh: bila $n=3$, maka $q=111$ adalah kelipatan dari $3$ yang semua digitnya adalah $1$.)

@ID: OW19D2P1U
@BIDANG: Aljabar Linear

Misalkan $V$ adalah ruang hasil kali dalam dan $\{u_1,u_2,\ldots,u_n\}$ adalah basis ortonormal $V$. Misalkan pula $\{v_1,v_2,\ldots,v_{n-1}\}\subset\{u_1,u_2,\ldots,u_n\}$.

Tentukan banyaknya vektor yang normanya $1$ dan ortogonal terhadap semua vektor $v_1,v_2,\ldots,v_{n-1}$.

@ID: OW19D2P2U
@BIDANG: Aljabar Linear

Misalkan $W$ adalah ruang vektor bagian dari ruang vektor $M_n(\mathbb R)$ yang memenuhi

$$
\operatorname{tr}(AB)=0
$$

untuk setiap $A,B\in W$.

Tentukan bilangan bulat terkecil $k$ sedemikian sehingga

$$
\dim W\le k.
$$

($M_n(\mathbb R)$ menyatakan himpunan semua matriks real berukuran $n\times n$.)

@ID: OW19D2P3U
@BIDANG: Analisis Kompleks

Diberikan bilangan kompleks $z_1,z_2,\ldots,z_n$ sehingga

$$
|z_1|=|z_2|=\cdots=|z_n|>0.
$$

Buktikan bahwa

$$
\operatorname{Re}\left(\sum_{j=1}^n\sum_{k=1}^n z_j\overline{z_k}\right)=0
$$

jika dan hanya jika

$$
\sum_{k=1}^n z_k=0.
$$

@ID: OW19D2P4U
@BIDANG: Kombinatorika

Diberikan bilangan bulat $n\ge4$. Tuliskan argumentasi kombinatorial untuk memperlihatkan bahwa

$$
\sum_{k=1}^n
k
\binom{n-4}{\,n-k\,}
\binom{n+4}{k}
=
(n+4)\binom{2n-1}{n-1}.
$$